% ============================================================
% vacuum_twist_field_equations.tex
% STATUS: STUB -- see notes/vacuum_twist_field_equations.md
% ============================================================

% This section derives the unified phase field equation and shows
% that gravity and electromagnetism are two types of vacuum deformation:
%   Gravity: div(phi) != 0  -- scalar, always attractive
%   EM:     curl(phi) != 0  -- vector, two charge signs

% 8.1 The Three Regimes of delta_phi
% ---------------------------------------------------------
% Rest mass:      delta_phi uniform -> symmetric Zitterbewegung
% Momentum:       constant phase gradient -> asymmetric bias -> drift
% Acceleration:   changing phase gradient -> changing bias -> steering
% These are Newton's three laws, derived from phase geometry.

% 8.2 Gravity as Divergence of the Phase Field
% ---------------------------------------------------------
% Clock density well creates div(phi) != 0.
% Phase flows toward the source (always attractive).
% Differential Zitterbewegung across wavefront steers packet.
% Schwarzschild metric emerges in continuum limit (Section 7).

% 8.3 Electromagnetism as Curl of the Phase Field
% ---------------------------------------------------------
% Vacuum twist creates curl(phi) != 0.
% Phase winds around the source (can attract or repel).
% Two winding directions = two charge signs.
% Maxwell's equations emerge in continuum limit.
% Connection to photon chirality: photons ARE the curl excitation.

% 8.4 The Unified Phase Field Equation
% ---------------------------------------------------------
% Box(phi) = -sin^2(phi/2)*phi + J_em + J_grav
% where:
%   Box = d^2/dt^2 - c^2 * Laplacian  (d'Alembertian)
%   J_em  = epsilon_munu * d^nu F^mu  (EM source)
%   J_grav = kappa * Laplacian(rho_clock)  (gravitational source)
%
% Special cases:
%   phi -> 0 (massless): wave equation -- photon
%   No sources:          Klein-Gordon equation
%   EM source only:      Maxwell's equations in Lorenz gauge
%   Grav source only:    linearized Einstein (weak field)

% 8.5 Acceleration as Differential Zitterbewegung
% ---------------------------------------------------------
% grad(delta_phi) across wavefront -> differential p_stay
% -> differential residence time -> packet steers
% No force vector required.
% Equivalence principle: gravitational and inertial mass
%   are the same p_stay = sin^2(delta_phi/2).

% 8.6 The Fine Structure Constant (Revised)
% ---------------------------------------------------------
% alpha ~ ratio of EM twist amplitude to lattice spacing
% The curl must be weak enough not to break bipartite vacuum structure.
% Stability threshold gives alpha ~ 1/137 geometrically.
% Cleaner derivation than v1.0 Gambler's Ruin -- to be completed
% in exp_08_vacuum_twist.py.

% 8.7 Solitons and Stable Particles
% ---------------------------------------------------------
% The nonlinear equation Box(phi) = -sin^2(phi/2)*phi
% may have stable soliton solutions.
% These solitons would be stable massive particles.
% The particle spectrum could emerge from soliton topology.
% Connection to topological quantum field theory.

\textbf{[STUB -- see notes/vacuum\_twist\_field\_equations.md]}
